\chapter*{Considerações Finais}\addcontentsline{toc}{chapter}{Considerações Finais}
Formou-se, no meio digital, uma indústria de poderio colossal fundamentada na coleta de dados pessoais.
Por meio deles, pôde criar modelos computacionais que sintetizam o indivíduo e, por meio disso, racionalizam sua existência a fim de otimizar a obtenção de lucro por parte dessas empresas.
Essa coleta ocorre às escuras, frequentemente mal-comunicada para os usuários desses serviços, que vivem cercados, em todos os lados de sua vida digital, por essa coleta.
Baseando-se em uma narrativa de libertação individual e oposição ao Estado, essas empresas se consolidaram, no auge do neoliberalismo pós-Guerra Fria, como monopólios centralizadores da experiência digital.
Passaram, por seus métodos indiretos de vigilância, a mediar a grande maioria de todo o tráfego pela Internet, de modo que hoje é quase inevitável entregá-los algum dado ao navegar.

Por meio de algoritmos de aprendizado de máquina sucessivamente mais refinados, consolidaram seu poder: esses algoritmos são o meio pelo qual o potencial dos dados é realizado, consumado, e transformado de um conceito abstrato em uma medida efetiva.
Em tempos recentes, esses métodos de processamento melhoraram, tanto em teoria quanto em potencial computacional; hoje, são necessários \textit{clusters} gigantescos de servidores, rodando algoritmos extremamente complexos, com milhões de parâmetros, para que esses serviços se sustentem.

Essas empresas constituíram-se também sobre uma base de obscurantismo e falta de transparência, coadjuvada pela natureza de caixa-preta desses modelos, mas que é, de fato, completamente arbitrária.
Essa maneira de produzir o espaço digital é uma imposição das indústrias do Vale do Silício, difundida ideologicamente pelas mesmas, de modo que conseguiu neutralizar muitas das correntes de pensamento opostas a ela, abarcando dentro das plataformas do Vale do Silício discursos dialeticamente opostos.

Há, na sociedade moderna, um problema: os maiores negócios da indústria da Internet fundamentam-se na expropriação e na ambiguidade, no monopólio econômico e político, e na alienação da consciência e do pensamento do indivíduo, delegado a algoritmos sintéticos em que existem, por se tratarem de aproximações, múltiplos vieses.

Essa monografia buscou iluminar as fundações trevosas sobre as quais a Internet é construída, ainda que a totalidade dessas páginas seja insuficiente para tal propósito: faltaram descrições mais aprofundadas em todas as dimensões: tanto das técnicas de coleta, quanto da matemática dos algoritmos, quanto da extensão política dessa indústria.
Apesar disso, espera-se que ela tenha sido satisfatória para o leitor, como um panorama do meio digital centrado na coleta de dados.
